\documentclass[12pt, a4paper]{article}
\usepackage[T2A]{fontenc}
\usepackage[utf8]{inputenc}
\usepackage[russian]{babel}
\usepackage{amsmath, amssymb, amsthm}
\usepackage{geometry}
\usepackage{bm}
\usepackage{arydshln} % для пунктирных линий в матрицах
\geometry{left=2.5cm, right=2cm, top=2cm, bottom=2cm}

\theoremstyle{definition}
\newtheorem*{theory}{Теория}
\newtheorem*{geom}{Геометрическое пояснение}
\newtheorem*{alg}{Алгебраическое пояснение}
\newtheorem{theorem}{Теорема}
\newtheorem*{remark}{Замечание}

\begin{document}

% Единый красивый заголовок
\begin{center}
    \Large\bfseries Задача №11: Анализ систем и аппроксимация матриц через SVD \\
\end{center}

\vspace{0.5cm}

\section{Теоретическая часть}

\begin{theory}
Для любой матрицы $A \in \mathbb{R}^{m \times n}$ существует разложение (сингулярные числа определены единственным образом):
\[
A = U \Sigma V^T,
\]
где:
\begin{itemize}
    \item $U \in \mathbb{R}^{m \times m}$, ортогональная матрица: $U^T U = I_m$. Столбцы $U = (u_1, \dots, u_m)$ --- левые сингулярные векторы.
    \item $V \in \mathbb{R}^{n \times n}$, ортогональная матрица: $V^T V = I_n$. Столбцы $V = (v_1, \dots, v_n)$ --- правые сингулярные векторы.
    \item $\Sigma \in \mathbb{R}^{m \times n}$, «диагональная» матрица вида $\Sigma = \begin{pmatrix} \sigma_1 & & 0 \\ & \ddots & \\ 0 & & \sigma_k \\ \hdashline & 0 & \end{pmatrix}$, где $\sigma_1 \ge \sigma_2 \ge \dots \ge \sigma_k > 0$ --- сингулярные числа. Также обозначают $\Sigma_n = \text{diag}(\sigma_1, \dots, \sigma_n)$.
\end{itemize}
\end{theory}

\begin{geom}
Отображение $A \colon \mathbb{R}^n \to \mathbb{R}^m$ переводит единичную сферу в пространстве $\mathbb{R}^n$ в эллипсоид в пространстве $\mathbb{R}^m$. Полуоси этого эллипсоида направлены вдоль левых сингулярных векторов $u_1, u_2, \dots, u_m$, а их длины равны соответствующим сингулярным числам $\sigma_1, \sigma_2, \dots, \sigma_n$. Остальные направления дополняются до ортогонального базиса.
\end{geom}

\begin{alg}
Для любого вектора $x \in \mathbb{R}^n$ его можно разложить по базису из правых сингулярных векторов:
\[
x = \sum_{i=1}^n c_i v_i \quad (\text{так как } V \text{ --- базис}).
\]
Тогда действие матрицы $A$ на вектор $x$ можно записать как:
\[
Ax = U \Sigma V^T x = (u_1, \dots, u_m) \Sigma V^T \left( \sum_{i=1}^n c_i v_i \right).
\]
Поскольку $V^T V = I_n$, имеем $V^T \left( \sum c_i v_i \right) = (c_1, \dots, c_n)^T$. Следовательно:
\[
Ax = \sum_{i=1}^n c_i \sigma_i u_i.
\]
\end{alg}

\hrule\vspace{0.5cm}

\begin{theorem}[О решении задачи наименьших квадратов, ранг полный]
Рассмотрим переопределенную систему $Ax = b$ (задача наименьших квадратов --- ЗНК). Пусть $\operatorname{rank}(A) = n$ (полный столбцовый ранг).
Тогда нормальное псевдорешение имеет вид:
\[
x = V \Sigma^{-1} U^T b.
\]
\end{theorem}
\begin{proof}
Задача наименьших квадратов сводится к решению нормального уравнения:
\[
A^T A x = A^T b \implies x = (A^T A)^{-1} A^T b.
\]
Подставим SVD-разложение $A = U \Sigma V^T$:
\begin{align*}
x &= (V \Sigma U^T U \Sigma V^T)^{-1} V \Sigma U^T b = \\
  &= (V \Sigma^2 V^T)^{-1} V \Sigma U^T b = \\
  &= V \Sigma^{-2} V^T V \Sigma U^T b = \\
  &= V \Sigma^{-1} U^T b.
\end{align*}
\end{proof}

\begin{theorem}[О решении ЗНК, ранг неполный]
Пусть $\operatorname{rank}(A) = k < n$. Тогда сингулярное разложение имеет блочный вид:
\[
A = U \begin{pmatrix} \Sigma_k & 0 \\ 0 & 0 \end{pmatrix} V^T.
\]
Тогда множество решений ЗНК имеет вид: $x = V_1 \Sigma_k^{-1} U_1^T b + V_2 z$, где $z \in \mathbb{R}^{n-k}$ --- произвольный вектор.
\end{theorem}
\begin{proof}
Минимизируем невязку:
\[
\|b - Ax\|_2 \to \inf_x \iff \|b - U \Sigma V^T x\|_2 \to \inf_x.
\]
Умножим выражение под нормой на ортогональную матрицу $U^T$ (умножение на ортогональную матрицу сохраняет евклидову норму, $\|U^T y\|_2 = \|y\|_2$):
\[
\|U^T b - \Sigma V^T x\|_2 \to \inf_x.
\]
Разобьем матрицы $U$ и $V$ на блоки, соответствующие ненулевым сингулярным числам: $U = (U_1, U_2)$ и $V = (V_1, V_2)$. Тогда $U^T b = \begin{pmatrix} U_1^T b \\ U_2^T b \end{pmatrix}$. Получаем:
\[
\left\| \begin{pmatrix} U_1^T b \\ U_2^T b \end{pmatrix} - \begin{pmatrix} \Sigma_k V_1^T x \\ 0 \end{pmatrix} \right\|_2 \to \inf_x \iff \left\| \begin{pmatrix} U_1^T b - \Sigma_k V_1^T x \\ U_2^T b \end{pmatrix} \right\|_2 \to \inf_x.
\]
Нижний блок $U_2^T b$ не зависит от $x$. Минимум достигается, когда верхний блок равен нулю. Обозначим $V^T x = \begin{pmatrix} V_1^T x \\ V_2^T x \end{pmatrix} = \begin{pmatrix} y_1 \\ y_2 \end{pmatrix}$. 
Решение для первой компоненты:
\[
y_1 = \Sigma_k^{-1} U_1^T b.
\]
Тогда искомый вектор $x$ восстанавливается как $x = V y = (V_1, V_2) \begin{pmatrix} y_1 \\ y_2 \end{pmatrix}$:
\[
x = V_1 y_1 + V_2 y_2 = V_1 \Sigma_k^{-1} U_1^T b + V_2 y_2, \quad \forall y_2 \in \mathbb{R}^{n-k}.
\]
\end{proof}

\begin{remark}
Если нас интересует вектор минимальной длины (нормальное псевдорешение), то $x_{\min} = V_1 \Sigma_k^{-1} U_1^T b$, то есть $y_2 = 0$.
\end{remark}
\begin{proof}
Поскольку столбцы матрицы $V$ ортонормальны, имеем ортогональное разложение:
\[
x = \sum_{i=1}^k c_i v_i + \sum_{j=k+1}^n \beta_j v_j, \quad \text{где } \beta_j \text{ --- компоненты } y_2.
\]
Длина вектора $\|x\|_2^2 = \sum c_i^2 + \sum \beta_j^2$. Мы хотим минимизировать длину. Так как коэффициенты $c_i$ (или $y_1$) жестко фиксированы для обеспечения минимума невязки, а $\beta_j$ на минимизацию невязки не влияют, то для минимизации длины самого вектора $x$ необходимо положить $\beta_j = 0$, то есть $y_2 = 0$.
\end{proof}

\hrule\vspace{0.5cm}

\begin{theorem}[Теорема Эккарта --- Янга --- Мирского]
Пусть $A \in \operatorname{Mat}_{m \times n}(\mathbb{R})$ и $\operatorname{rank}(A) = n$. Тогда наилучшая аппроксимация матрицы $A$ матрицей меньшего ранга $k$ в норме $\|\cdot\|_2$:
\[
\inf_{\substack{B \in \operatorname{Mat}_{m \times n}(\mathbb{R}) \\ \operatorname{rank}(B) \le k}} \|A - B\|_2
\]
достигается (и притом единственно) на матрице усеченного SVD-разложения:
\[
A_k = U \begin{pmatrix} \Sigma_k & 0 \\ 0 & 0 \end{pmatrix} V^T.
\]
\end{theorem}
\begin{proof}\mbox{}
\begin{enumerate}
    \item Очевидно, что $A_k$ имеет ранг $k$. Рассмотрим разность:
    \[
    A - A_k = U \begin{pmatrix} 0 & 0 \\ 0 & \Sigma_{n-k} \end{pmatrix} V^T.
    \]
    Покажем, что $\|A - A_k\|_2 = \sigma_{k+1}$ (где $\sigma_{k+1}$ --- максимальное сингулярное число в блоке $\Sigma_{n-k}$, так как на диагонали стоят обнуленные $\sigma_1 \dots \sigma_k$). \\
    Действительно, норма матрицы равна максимальному сингулярному числу, поэтому $\|A - A_k\|_2 \le \sigma_{k+1}$. С другой стороны:
    \[
    \|A - A_k\|_2 \ge \frac{\|(A - A_k)v_{k+1}\|_2}{\|v_{k+1}\|_2} = \|\sigma_{k+1}u_{k+1}\|_2 = \sigma_{k+1}.
    \]
    Следовательно, $\|A - A_k\|_2 = \sigma_{k+1}$.

    \item Пусть теперь $B$ --- любая матрица такая, что $\operatorname{rank}(B) = k$. 
    Рассмотрим ядро матрицы $B$, то есть подпространство $\ker B$. По теореме о ранге и дефекте, $\dim(\ker B) = n - k$. 
    Рассмотрим также подпространство $\operatorname{span}\langle v_1, \dots, v_{k+1} \rangle$, его размерность равна $k+1$.
    Поскольку сумма их размерностей $(n - k) + (k + 1) = n + 1 > n$, эти два подпространства пересекаются нетривиально:
    \[
    \exists \xi \in \ker B \cap \operatorname{span}\langle v_1, \dots, v_{k+1} \rangle, \quad \xi \neq 0.
    \]
    Не умаляя общности, нормируем этот вектор: $\|\xi\|_2 = 1$. Разложим его по базису:
    $\xi = \sum_{i=1}^{k+1} \alpha_i v_i$, причем $\sum \alpha_i^2 = 1$. Так как $\xi \in \ker B$, то $B\xi = 0$.
    
    Тогда оценим норму разности операторов на этом векторе $\xi$:
    \[
    (A - B)\xi = A\xi - 0 = A\left( \sum_{i=1}^{k+1} \alpha_i v_i \right) = \sum_{i=1}^{k+1} \alpha_i \sigma_i u_i.
    \]
    Норма матрицы не меньше нормы ее действия на единичный вектор:
    \[
    \|A - B\|_2 \ge \|(A - B)\xi\|_2 = \left\| \sum_{i=1}^{k+1} \alpha_i \sigma_i u_i \right\|_2 = \sqrt{\sum_{i=1}^{k+1} (\alpha_i \sigma_i)^2}.
    \]
    Поскольку сингулярные числа упорядочены по убыванию, $\sigma_i \ge \sigma_{k+1}$ для всех $i \le k+1$:
    \[
    \|A - B\|_2 \ge \sqrt{\sum_{i=1}^{k+1} \alpha_i^2 \sigma_{k+1}^2} = \sigma_{k+1} \sqrt{\sum_{i=1}^{k+1} \alpha_i^2} = \sigma_{k+1} \cdot 1 = \sigma_{k+1}.
    \]
    Таким образом, для любой матрицы $B$ ранга $k$ ошибка аппроксимации $\|A - B\|_2 \ge \sigma_{k+1}$. Минимум достигается на $A_k$.
\end{enumerate}
\end{proof}

\vspace{0.5cm}

\section{Численные эксперименты}

В рамках работы была реализована программа на языке Python (с использованием библиотек \texttt{numpy} и \texttt{scipy.linalg.svd}), тестирующая метод сингулярного разложения на трех типах матриц.

\subsection{Случай 1: Плохо обусловленная верхнетреугольная матрица}
Рассматривается матрица $A \in \mathbb{R}^{100 \times 100}$ с единицами на диагонали и $(-1)$ выше неё. 
\begin{itemize}
    \item \textbf{Проблема:} Матрица формально невырожденная, но спектр сингулярных чисел $\sigma_i$ убывает крайне быстро (от $\approx 5.6$ до $\approx 10^{-3}$).
    \item \textbf{Результат:} В точной арифметике система $Ax=b$ разрешима. Однако на практике число обусловленности $\kappa(A) = \sigma_1 / \sigma_n$ велико. Использование усеченного SVD (с отсечением сингулярных чисел меньше заданного $\varepsilon=10^{-12}$) позволяет стабильно найти нормальное псевдорешение.
\end{itemize}

\subsection{Случай 2: Вырожденная трехдиагональная матрица}
Рассматривается матрица $A$ со значениями $(2, -1, -1)$ на диагоналях и модифицированными граничными условиями (значения $-2$ в углах).
\begin{itemize}
    \item \textbf{Особенность:} Матрица вырождена, $\operatorname{rank}(A) = 99 < 100$. В процессе вычислений одно из сингулярных чисел определяется как машинный ноль ($\approx 10^{-16}$). Стандартный метод Гаусса (или \texttt{np.linalg.solve}) к такой системе неприменим.
    \item \textbf{Метод решения:} Метод SVD позволил выделить вектор $v_n$ из матрицы $V$ (соответствующий $\sigma_n \approx 0$), который задает базис одномерного ядра матрицы $A$. Общее решение восстанавливается в виде $x = x_{\min} + C \cdot v_n$. Расчеты показали, что ошибка восстановления точного вектора таким образом близка к машинному нулю.
\end{itemize}

\subsection{Случай 3: Низкоранговая аппроксимация (Синус-матрица)}
Рассматривается матрица размерности $5 \times 5$, элементы которой заданы как $A_{i,j} = \sin(\pi x_i) \sin(\pi y_j)$, где $x_i, y_j$ --- узлы равномерной сетки.
\begin{itemize}
    \item \textbf{Анализ:} Поскольку переменные $x$ и $y$ сепарабельны (каждый элемент есть произведение функции от $i$ на функцию от $j$), строки матрицы линейно зависимы. Теоретический ранг такой матрицы равен 1. 
    \item \textbf{Результат:} Анализ спектра сингулярных значений подтверждает теорию: только $\sigma_1$ отлично от нуля (с точностью до ошибок округления), что позволяет сжимать подобные данные без потерь, оставляя только первую компоненту усеченного SVD.
\end{itemize}

\end{document}